%==========================================================
%  CHAPITRE 4 — SUITES NUMÉRIQUES
%==========================================================

\chapter{Suites Num\'{e}riques}

\begin{rappel}[Vocabulaire de base]
  \begin{itemize}
    \item Une \textbf{suite num\'{e}rique} est une fonction $u : \N \to \R$,
          not\'{e}e $(u_n)_{n \in \N}$.
    \item On peut d\'{e}finir une suite par une \textbf{formule explicite}
          $u_n = f(n)$ ou par une \textbf{relation de r\'{e}currence}
          $u_{n+1} = f(u_n)$.
    \item Le \textbf{terme g\'{e}n\'{e}ral} $u_n$ est la valeur de la suite au rang $n$.
  \end{itemize}
\end{rappel}

%----------------------------------------------------------
\section{Suites arithm\'{e}tiques et g\'{e}om\'{e}triques}
%----------------------------------------------------------

\begin{definition}[Suite arithm\'{e}tique]
  Une suite $(u_n)$ est \textbf{arithm\'{e}tique} de raison $r$ si :
  \[
    \forall n \in \N, \quad u_{n+1} = u_n + r
  \]
\end{definition}

\begin{definition}[Suite g\'{e}om\'{e}trique]
  Une suite $(u_n)$ est \textbf{g\'{e}om\'{e}trique} de raison $q$ si :
  \[
    \forall n \in \N, \quad u_{n+1} = q \cdot u_n \qquad (q \neq 0,\ u_0 \neq 0)
  \]
\end{definition}
\newpage 
\begin{propriete}[Tableau comparatif]
  \bigskip
  \renewcommand{\arraystretch}{2.2}
  \begin{tabularx}{\linewidth}{|>{\bfseries}p{3.2cm}|>{\centering\arraybackslash}X|>{\centering\arraybackslash}X|}
    \hline
    \rowcolor{white}
      & \textbf{\color{navy}Suite arithm\'{e}tique}
      & \textbf{\color{forest}Suite g\'{e}om\'{e}trique} \\
    \hline\hline
    D\'{e}finition
      & $u_{n+1} = u_n + r$
      & $u_{n+1} = q \cdot u_n$ \\
    \hline
    Terme g\'{e}n\'{e}ral
      & $u_n = u_p + (n-p)\,r$
      & $u_n = u_p \cdot q^{n-p}$ \\
    \hline
    Somme de termes cons\'{e}cutifs
      & $S = \dfrac{(n-p+1)(u_p + u_n)}{2}$
      & $S = u_p \cdot \dfrac{1 - q^{n-p+1}}{1-q}$ \; ($q \neq 1$) \\
    \hline
    $a$, $b$, $c$ cons\'{e}cutifs
      & $2b = a + c$
      & $b^2 = a \cdot c$ \\
    \hline
  \end{tabularx}
\end{propriete}

\begin{remarque}
  Pour la somme g\'{e}om\'{e}trique, si $q = 1$ alors $S = (n-p+1)\,u_p$.

  \medskip
  Formule utile : $1 + q + q^2 + \cdots + q^n = \dfrac{1 - q^{n+1}}{1-q}$
  pour $q \neq 1$.
\end{remarque}

\begin{exemple}[R\'{e}solu : suite arithm\'{e}tique]
  $(u_n)$ arithm\'{e}tique avec $u_0 = 3$ et $r = 5$.

  \medskip
  \textbf{Terme g\'{e}n\'{e}ral :} $u_n = 3 + 5n$.

  \medskip
  \textbf{Somme :} $u_0 + u_1 + \cdots + u_{10}
  = \dfrac{11 \times (u_0 + u_{10})}{2}
  = \dfrac{11 \times (3 + 53)}{2} = \dfrac{11 \times 56}{2} = 308$.
\end{exemple}

\begin{exemple}[R\'{e}solu : suite g\'{e}om\'{e}trique]
  $(u_n)$ g\'{e}om\'{e}trique avec $u_0 = 2$ et $q = 3$.

  \medskip
  \textbf{Terme g\'{e}n\'{e}ral :} $u_n = 2 \times 3^n$.

  \medskip
  \textbf{Somme :} $u_0 + u_1 + \cdots + u_4
  = 2 \times \dfrac{1 - 3^5}{1 - 3}
  = 2 \times \dfrac{1 - 243}{-2} = 242$.
\end{exemple}
\newpage
%----------------------------------------------------------
\section{Monotonie et bornitude}
%----------------------------------------------------------

\begin{definition}[Monotonie d'une suite]
  Soit $(u_n)$ une suite num\'{e}rique.
  \begin{itemize}
    \item $(u_n)$ est \textbf{croissante} si :
          $\forall n \in \N,\quad u_{n+1} \geq u_n$
    \item $(u_n)$ est \textbf{d\'{e}croissante} si :
          $\forall n \in \N,\quad u_{n+1} \leq u_n$
    \item $(u_n)$ est \textbf{constante} si :
          $\forall n \in \N,\quad u_{n+1} = u_n$
  \end{itemize}
\end{definition}

\begin{methode}[\'{E}tudier la monotonie d'une suite]
  On peut \'{e}tudier le signe de $u_{n+1} - u_n$, ou,
  si les termes sont positifs, comparer $\dfrac{u_{n+1}}{u_n}$ \`{a} $1$.

  \medskip
  \begin{center}
  \renewcommand{\arraystretch}{1.8}
  \begin{tabular}{|c|c|}
    \hline
    \textbf{\color{navy}Condition}
      & \textbf{\color{navy}Conclusion} \\
    \hline\hline
    $u_{n+1} - u_n > 0$
      & $(u_n)$ \textbf{croissante} \\
    \hline
    $u_{n+1} - u_n < 0$
      & $(u_n)$ \textbf{d\'{e}croissante} \\
    \hline
    $\dfrac{u_{n+1}}{u_n} > 1$ \; ($u_n > 0$)
      & $(u_n)$ \textbf{croissante} \\
    \hline
    $\dfrac{u_{n+1}}{u_n} < 1$ \; ($u_n > 0$)
      & $(u_n)$ \textbf{d\'{e}croissante} \\
    \hline
  \end{tabular}
  \end{center}
\end{methode}

\begin{definition}[Suite born\'{e}e]
  Soit $(u_n)$ une suite num\'{e}rique.
  \begin{itemize}
    \item $(u_n)$ est \textbf{major\'{e}e} par $M$ si :
          $\forall n \in \N,\quad u_n \leq M$
    \item $(u_n)$ est \textbf{minor\'{e}e} par $m$ si :
          $\forall n \in \N,\quad u_n \geq m$
    \item $(u_n)$ est \textbf{born\'{e}e} si elle est \`{a} la fois
          major\'{e}e et minor\'{e}e.
  \end{itemize}
\end{definition}

\begin{exemple}[R\'{e}solu : monotonie et bornitude]
  $u_n = \dfrac{n}{n+1}$.

  \medskip
  $u_{n+1} - u_n
  = \dfrac{n+1}{n+2} - \dfrac{n}{n+1}
  = \dfrac{(n+1)^2 - n(n+2)}{(n+2)(n+1)}
  = \dfrac{1}{(n+2)(n+1)} > 0$.

  Donc $(u_n)$ est \textbf{croissante}.

  De plus $0 \leq u_n < 1$ pour tout $n$ :
  $(u_n)$ est \textbf{born\'{e}e} (minor\'{e}e par $0$, major\'{e}e par $1$).
\end{exemple}

%----------------------------------------------------------
\section{Limite d'une suite}
%----------------------------------------------------------

\begin{definition}[Convergence et divergence]
  \begin{itemize}
    \item $(u_n)$ \textbf{converge} vers $\ell \in \R$ si $u_n$
          peut \^{e}tre rendu aussi proche de $\ell$ que l'on veut
          pour $n$ suffisamment grand. On \'{e}crit :
          $\displaystyle\lim_{n \to +\infty} u_n = \ell$
    \item $(u_n)$ \textbf{diverge vers $+\infty$} si $u_n$ d\'{e}passe
          tout r\'{e}el pour $n$ grand.
  \end{itemize}
\end{definition}

\begin{propriete}[Limite de $n^\alpha$ avec $\alpha \in \R^*$]
  \[
    \lim_{n \to +\infty} n^\alpha = +\infty \quad (\alpha > 0)
    \qquad\qquad
    \lim_{n \to +\infty} n^\alpha = 0 \quad (\alpha < 0)
  \]
\end{propriete}

\begin{propriete}[Limite de $q^n$ selon la valeur de $q$]
  \bigskip
  \begin{center}
  \renewcommand{\arraystretch}{1.8}
  \begin{tabular}{|c|c|c|c|c|}
    \hline
    \textbf{\color{navy}$q < -1$}
      & \textbf{\color{navy}$q = -1$}
      & \textbf{\color{navy}$-1 < q < 1$}
      & \textbf{\color{navy}$q = 1$}
      & \textbf{\color{navy}$q > 1$} \\
    \hline\hline
    Pas de limite & Pas de limite & $0$ & $1$ & $+\infty$ \\
    \hline
  \end{tabular}
  \end{center}
\end{propriete}

\begin{propriete}[Th\'{e}or\`{e}me des gendarmes pour les suites]
  Si $\forall n \in \N : \; U_n \leq u_n \leq V_n$
  et si $\displaystyle\lim_{n\to+\infty} U_n
       = \lim_{n\to+\infty} V_n = \ell$,
  alors :
  \[
    \lim_{n \to +\infty} u_n = \ell
  \]
\end{propriete}

\begin{propriete}[Th\'{e}or\`{e}me de comparaison]
  \begin{itemize}
    \item Si $u_n \leq v_n$ et $\displaystyle\lim u_n = +\infty$,
          alors $\displaystyle\lim v_n = +\infty$.
    \item Si $u_n \leq v_n$ pour tout $n$ et $(u_n)$, $(v_n)$
          convergent, alors $\displaystyle\lim u_n \leq \lim v_n$.
  \end{itemize}
\end{propriete}
\newpage

%----------------------------------------------------------
\section{Crit\`{e}res de convergence}
%----------------------------------------------------------

\begin{theoreme}[Suites monotones et born\'{e}es]
  \begin{itemize}
    \item Toute suite \textbf{croissante et major\'{e}e} est \textbf{convergente}.
    \item Toute suite \textbf{d\'{e}croissante et minor\'{e}e} est \textbf{convergente}.
  \end{itemize}
\end{theoreme}

\begin{attention}
  Ces th\'{e}or\`{e}mes affirment l'existence de la limite,
  mais \textbf{ne donnent pas sa valeur}.
  Pour trouver $\ell$, on utilise les m\'{e}thodes des sections suivantes.
\end{attention}

%----------------------------------------------------------
\section{Suite de type \texorpdfstring{$v_n = f(u_n)$}{v\_n = f(u\_n)}}
%----------------------------------------------------------

\begin{theoreme}[Suite compos\'{e}e]
  Si $(u_n)$ est convergente de limite $\ell$ et si $f$ est
  \textbf{continue en $\ell$}, alors la suite $(v_n)$ d\'{e}finie par
  $v_n = f(u_n)$ est convergente et :
  \[
    \lim_{n \to +\infty} v_n = f(\ell)
  \]
\end{theoreme}

\begin{exemple}[R\'{e}solu]
  $u_n = \dfrac{1}{n}$ (converge vers $0$).
  Soit $v_n = e^{u_n} = e^{1/n}$.

  $f(x) = e^x$ est continue en $0$, donc :
  $\displaystyle\lim_{n\to+\infty} v_n = e^0 = 1$.
\end{exemple}

%----------------------------------------------------------
\section{Suite de type \texorpdfstring{$u_{n+1} = f(u_n)$}{u\_{n+1} = f(u\_n)}}
%----------------------------------------------------------

\begin{theoreme}[Suite r\'{e}currente]
  Soit $(u_n)$ d\'{e}finie par $u_0 = a$ et $u_{n+1} = f(u_n)$.
  Si :
  \begin{itemize}
    \item $f$ est \textbf{continue} sur un intervalle $I$
    \item $f(I) \subset I$ \quad (stabilit\'{e})
    \item $a \in I$
    \item $(u_n)$ est \textbf{convergente} de limite $\ell$
  \end{itemize}
  alors $\ell$ est \textbf{solution de l'\'{e}quation} $f(x) = x$
  (point fixe de $f$).
\end{theoreme}

\begin{methode}[\'{E}tudier une suite r\'{e}currente $u_{n+1} = f(u_n)$]
  \begin{enumerate}
    \item \textbf{Montrer} que $(u_n)$ est dans un intervalle $I$
          stable par $f$ (par r\'{e}currence).
    \item \textbf{Montrer} la monotonie (croissante ou d\'{e}croissante).
    \item \textbf{Conclure} \`{a} la convergence (born\'{e}e + monotone).
    \item \textbf{Trouver} la limite $\ell$ en r\'{e}solvant $f(\ell) = \ell$.
  \end{enumerate}
\end{methode}

\begin{exemple}[R\'{e}solu : suite r\'{e}currente]
  $u_0 = 0$ et $u_{n+1} = \dfrac{u_n + 2}{2}$.

  \medskip
  \textbf{\'{E}tape 1 :} Montrons par r\'{e}currence que $0 \leq u_n \leq 2$.

  $u_0 = 0 \in [0,2]$. Si $u_n \in [0,2]$ alors
  $u_{n+1} = \dfrac{u_n+2}{2} \in [1,2] \subset [0,2]$. $\checkmark$

  \medskip
  \textbf{\'{E}tape 2 :}
  $u_{n+1} - u_n = \dfrac{u_n+2}{2} - u_n = \dfrac{2 - u_n}{2} \geq 0$
  car $u_n \leq 2$.
  Donc $(u_n)$ est \textbf{croissante}.

  \medskip
  \textbf{\'{E}tape 3 :} Croissante et major\'{e}e par $2$ $\Rightarrow$
  $(u_n)$ \textbf{converge}.

  \medskip
  \textbf{\'{E}tape 4 :} $\ell = f(\ell)$ :
  $\ell = \dfrac{\ell + 2}{2} \Rightarrow 2\ell = \ell + 2 \Rightarrow \ell = 2$.
\end{exemple}
\newpage
%----------------------------------------------------------
\section{Repr\'{e}sentation graphique d'une suite r\'{e}currente}
%----------------------------------------------------------
\begin{propriete}[Toile d'araign\'{e}e]
  La suite $u_{n+1} = f(u_n)$ se visualise graphiquement par
  la \textbf{m\'{e}thode de la toile d'araign\'{e}e} :
  \begin{enumerate}
    \item Tracer $\mathcal{C}_f$ (courbe de $f$) et la droite $y = x$.
    \item Partir du point $(u_0, 0)$ sur l'axe des abscisses.
    \item Monter verticalement jusqu'\`{a} $\mathcal{C}_f$ :
          on obtient $(u_0, f(u_0)) = (u_0, u_1)$.
    \item Aller horizontalement jusqu'\`{a} la droite $y=x$ :
          on obtient $(u_1, u_1)$.
    \item R\'{e}p\'{e}ter les \'{e}tapes 3--4.
  \end{enumerate}
\end{propriete}

\begin{exemple}[Toile d'araign\'{e}e : $u_{n+1} = \frac{u_n+2}{2}$, $u_0 = 0$]
  \begin{center}
  \begin{tikzpicture}
    \begin{axis}[
      axis lines  = center,
      xlabel      = {$x$},
      ylabel      = {$y$},
      xmin = -0.2, xmax = 2.8,
      ymin = -0.2, ymax = 2.8,
      xtick       = {0,0.5,1,1.5,2,2.5},
      ytick       = {0,0.5,1,1.5,2,2.5},
      grid        = major,
      grid style  = {gray!20},
      width = 10cm, height = 10cm,
    ]
      % Courbe f(x) = (x+2)/2
      \addplot[navy, very thick, domain=0:2.8]
        {(x+2)/2};
      \addlegendentry{$f(x)=\frac{x+2}{2}$}

      % Droite y = x
      \addplot[burgundy, dashed, thick, domain=0:2.8]
        {x};
      \addlegendentry{$y = x$}

      % Point fixe ℓ = 2
      \addplot[mark=*, mark size=4pt, burgundy]
        coordinates {(2,2)};
      \node[above left, burgundy, font=\small\bfseries]
        at (axis cs:2,2) {$\ell=2$};

      % Toile d'araignée : u0=0
      % (0,0) → (0, f(0)=1)
      \addplot[forest, thick]
        coordinates {(0,0)(0,1)};
      % (0,1) → (1,1)
      \addplot[forest, thick]
        coordinates {(0,1)(1,1)};
      % (1,1) → (1, f(1)=1.5)
      \addplot[forest, thick]
        coordinates {(1,1)(1,1.5)};
      % (1,1.5) → (1.5,1.5)
      \addplot[forest, thick]
        coordinates {(1,1.5)(1.5,1.5)};
      % (1.5,1.5) → (1.5, f(1.5)=1.75)
      \addplot[forest, thick]
        coordinates {(1.5,1.5)(1.5,1.75)};
      % (1.5,1.75) → (1.75,1.75)
      \addplot[forest, thick]
        coordinates {(1.5,1.75)(1.75,1.75)};
      % (1.75,1.75) → (1.75, f(1.75)=1.875)
      \addplot[forest, thick]
        coordinates {(1.75,1.75)(1.75,1.875)};
      % (1.75,1.875) → (1.875,1.875)
      \addplot[forest, thick]
        coordinates {(1.75,1.875)(1.875,1.875)};

      % Points u0, u1, u2, u3
      \addplot[mark=|, mark size=5pt, forest]
        coordinates {(0,0)(1,0)(1.5,0)(1.75,0)};
      \node[below, forest, font=\footnotesize] at (axis cs:0,0)    {$u_0$};
      \node[below, forest, font=\footnotesize] at (axis cs:1,0)    {$u_1$};
      \node[below, forest, font=\footnotesize] at (axis cs:1.5,0)  {$u_2$};
      \node[below, forest, font=\footnotesize] at (axis cs:1.75,0) {$u_3$};
    \end{axis}
  \end{tikzpicture}
  \end{center}
\end{exemple}

\newpage
%----------------------------------------------------------
\section*{Exercices du chapitre}
\addcontentsline{toc}{section}{Exercices : Suites num\'{e}riques}
%----------------------------------------------------------

\begin{exercice}[1 : Suite arithm\'{e}tique]
  $(u_n)$ est arithm\'{e}tique avec $u_2 = 7$ et $u_5 = 16$.
  \begin{enumerate}
    \item D\'{e}terminer la raison $r$ et le premier terme $u_0$.
    \item Exprimer $u_n$ en fonction de $n$.
    \item Calculer $S = u_0 + u_1 + \cdots + u_{20}$.
  \end{enumerate}
\end{exercice}

\begin{exercice}[2 : Suite g\'{e}om\'{e}trique]
  $(u_n)$ est g\'{e}om\'{e}trique avec $u_0 = 5$ et $u_3 = 40$.
  \begin{enumerate}
    \item D\'{e}terminer la raison $q$.
    \item Exprimer $u_n$ en fonction de $n$.
    \item Calculer $S = u_0 + u_1 + \cdots + u_7$.
  \end{enumerate}
\end{exercice}

\begin{exercice}[3 : Monotonie et bornitude]
  Pour chaque suite, \'{e}tudier la monotonie et la bornitude :
  \begin{enumerate}
    \item $u_n = \dfrac{2n+1}{n+2}$
    \item $u_n = 2^n - n$
    \item $u_n = \dfrac{(-1)^n}{n+1}$
  \end{enumerate}
\end{exercice}

\begin{exercice}[4 : Limite d'une suite]
  Calculer les limites suivantes :
  \begin{multicols}{2}
  \begin{enumerate}
    \item $\displaystyle\lim_{n\to+\infty} \frac{3n^2-1}{n^2+4}$
    \item $\displaystyle\lim_{n\to+\infty} \frac{n+1}{\sqrt{n}}$
    \item $\displaystyle\lim_{n\to+\infty} \left(\frac{2}{3}\right)^n$
    \item $\displaystyle\lim_{n\to+\infty} n \cdot \left(\frac{1}{2}\right)^n$
  \end{enumerate}
  \end{multicols}
\end{exercice}

\begin{exercice}[5 : Suite r\'{e}currente]
  Soit $(u_n)$ d\'{e}finie par $u_0 = 1$ et $u_{n+1} = \sqrt{u_n + 2}$.
  \begin{enumerate}
    \item Montrer par r\'{e}currence que $u_n \in [0, 2]$ pour tout $n$.
    \item Montrer que $(u_n)$ est croissante.
    \item Conclure que $(u_n)$ converge et d\'{e}terminer sa limite.
  \end{enumerate}
\end{exercice}

\begin{exercice}[6 : Toile d'araign\'{e}e]
  Soit $(u_n)$ d\'{e}finie par $u_0 = 3$ et $u_{n+1} = \dfrac{1}{2}u_n + 1$.
  \begin{enumerate}
    \item R\'{e}soudre $f(x) = x$ pour trouver le point fixe.
    \item Calculer $u_1$, $u_2$, $u_3$.
    \item Repr\'{e}senter la toile d'araign\'{e}e.
    \item D\'{e}terminer la limite de $(u_n)$.
  \end{enumerate}
\end{exercice}